\section{Milieux magnétiques}

    On peut distinguer expérimentalement dans la vie quotidienne différents types de magnétismes :
    \begin{itemize}
        \item \textbf{ferromagnétisme :} Un clou en fer est attiré par un aimant, et conserve une aimantation rémanente après retrait de l’aimant. 
        \item \textbf{paramagnétisme :} Un fourchette en aluminium est attirée plus faiblement par un aimant qu’un clou en fer, et il n’y a pas d’aimantation rémanente.
        \item \textbf{diamagnétisme :} Certains matériaux sont repoussés par un aimant, comme l’eau, de façon très faible.
    \end{itemize}
    Nous essayerons dans ce chapitre de comprendre l’origine microscopique de ces phénomènes.

    \subsection{Moment magnétique}

        L’origine microscopique du magnétisme provient pour la plus grande part des moments magnétiques des électrons dans les atomes. Ceux-ci sont proportionnels au moment angulaire orbital $\mathbf{L}$ et au moment angulaire de spin $\mathbf{S}$ électronique :
        \begin{align*}
            \mu_L &= - 1 \frac{\mu_B}{\hbar} \mathbf{L} \\
            \mu_S &= - 2 \frac{\mu_B}{\hbar} \mathbf{S} \\
        \end{align*}
        où le magnéton de Bohr est le « quantum » du moment magnétique $\mu_B = \frac{e\hbar}{2m} \approx 9 \times 10^{-24} \, \text{A.m}^2$.

        Pour calculer le champ magnétique créé par un moment magnétique, utilisons le fait expérimental que celui-ci est similaire à celui créé par une spire de rayon $a$ de surface $S= \pi a^2$ parcourue par un courant $I$ que l’on appelle un \textbf{dipôle magnétique}. Le potentiel vecteur résultant de la présence du dipôle magnétique est dans l’approximation dipolaire $(r \gg a)$ :
        \begin{align*}
            \mathbf{A}(\mathbf{r}) 
            &= \frac{\mu_0}{4\pi} \oint_{C} \frac{I d\mathbf{r}'}{\abs{\mathbf{r} - \mathbf{r}'}} \\
            &= \frac{\mu_0 I}{4 \pi} \int_{\Sigma} d\mathbf{S} \times \nabla_{r'} \left(\frac{1}{\abs{\mathbf{r} - \mathbf{r}'}}\right) \\
            &\approx \frac{\mu_0}{4\pi} \left(I \int_{\Sigma} d \mathbf{S}\right) \times \frac{\mathbf{e}_r}{r^2} \\ 
            &= \frac{\mu_0}{4\pi} \frac{\mu \times \mathbf{e}_r}{r^2}
        \end{align*}
        où $\mu = I \int_{\Sigma} d \mathbf{S}$ est le moment magnétique. Le champ magnétique vaut alors 
        \begin{align*}
            \mathbf{B}(\mathbf{r}) 
            &= \nabla \times \mathbf{A} \\
            &= \frac{\mu_0}{4\pi} \frac{3(\mathbf{e}_r \cdot \mu)\mathbf{e}_r - \mu}{r^3} \\
        \end{align*}

        